\section{Método de Revisão}

A Introdução situou os frameworks de apoio ao desenvolvimento de software com IA como uma camada que organiza o trabalho de quem usa um agente de desenvolvimento, e argumentou que compará-los exige examinar a função de cada mecanismo, e não métricas superficiais como número de comandos ou popularidade. Investigar esse objeto impõe uma restrição metodológica: os frameworks mais relevantes circulam fora de venues acadêmicos, em repositórios abertos, documentação oficial e kits comunitários, de modo que uma revisão restrita a artigos indexados deixaria de cobri-los. Por isso, esta revisão foi conduzida como uma revisão de literatura multivocal (Multivocal Literature Review, MLR), que combina, sob um único protocolo, o rigor da revisão sistemática de Kitchenham para a literatura formal e as diretrizes de Garousi, Felderer e Mäntylä para a literatura cinzenta \cite{kitchenham2009slr,garousi2019mlr}. A data de corte foi 26 de maio de 2026 e a janela de cobertura, de 2018 a 2026.

Os dois corpora cumprem papéis distintos e complementares. A literatura formal, composta por artigos e preprints recuperados em bases acadêmicas, fundamenta as dimensões de análise da taxonomia e estabelece o estado do conhecimento sobre agentes, especificação, validação e riscos. A literatura cinzenta, composta pelas fontes primárias dos frameworks de apoio, fornece os objetos efetivamente caracterizados por essas dimensões. O protocolo foi fixado antes da execução e define as perguntas de pesquisa, as fontes consultadas, os critérios de inclusão e exclusão, o esquema de extração e a estratégia de snowballing. A Figura~\ref{fig:metodologia} resume o fluxo de levantamento da literatura formal, da identificação ao relato, com as contagens reais de cada etapa.

\begin{figure}[h!]
\centering
\includegraphics[width=0.92\textwidth]{methodology_pt.pdf}
\caption{Fluxo da revisão de literatura multivocal aplicado à literatura formal, do protocolo Kitchenham ao relato, com as contagens reais de cada etapa.}
\label{fig:metodologia}
\end{figure}

A revisão é orientada pelas cinco perguntas de pesquisa apresentadas na Introdução (RQ1 a RQ5); a RQ5, sobre as lacunas que permanecem para comparar, avaliar e governar esses frameworks em projetos reais, é uma questão de nível de síntese, respondida transversalmente na Agenda de Pesquisa, e por isso a matriz de extração registra apenas RQ1 a RQ4 por estudo.

A busca seguiu duas estratégias. Para a literatura formal, consultaram-se de forma automatizada arXiv, ACM Digital Library, IEEE Xplore, Crossref, OpenAlex e Semantic Scholar, com Scopus e Web of Science acessados por metadados abertos via OpenAlex e o Google Scholar usado para verificação de citações. Os descritores foram organizados em quatro blocos conceituais (desenvolvimento assistido por IA, engenharia de software agêntica, especificação como artefato e organização por workflows), com strings booleanas adaptadas a cada base e documentadas no protocolo. Para a literatura cinzenta, em que a busca por string não se aplica, adotou-se busca dirigida às fontes primárias dos frameworks (documentação oficial, repositórios e blogs de produto), com critérios próprios de escopo e tração detalhados adiante; o esquema de extração, porém, foi o mesmo nas duas vertentes.

Os critérios de elegibilidade da literatura formal foram fixados no protocolo e estão resumidos na Tabela~\ref{tab:criterios}. A regra de decisão é agregada: uma fonte que não satisfaz o critério I1 é excluída; entre as fontes que satisfazem I1, são incluídas apenas as que satisfazem ao menos dois dos critérios I2, I3 e I4. Os motivos de exclusão foram registrados de forma operacional para cada fonte descartada, permitindo o relato auditável das razões de descarte em cada passada.

\begin{table}[h!]
\centering
\caption{Critérios de inclusão e exclusão aplicados na seleção dos estudos.}
\label{tab:criterios}
\small
\begin{tabularx}{\textwidth}{@{}p{0.46\textwidth}Y@{}}
\toprule
\textbf{Critérios de inclusão} & \textbf{Critérios de exclusão} \\
\midrule
I1: descreve framework, método, toolkit, plataforma, IDE, CLI ou kit que organiza o desenvolvimento com IA em mais de uma etapa. & E1: trata apenas de modelos de linguagem, APIs ou benchmarks de código, sem camada de processo. \\
I2: descreve produção ou consumo de artefatos persistentes (specs, planos, PRDs, tarefas, logs, evidências). & E2: descreve assistente de código genérico, sem artefatos persistentes, papéis, comandos ou validação própria. \\
I3: define papéis, agentes, personas, skills, workflows ou comandos de coordenação. & E3: discute LLMs para uma única tarefa isolada, sem proposta de framework ou processo. \\
I4: descreve integração com execução de código, testes, terminal, IDE, navegador ou repositórios. & E4: é puramente promocional, sem critérios, arquitetura ou evidência rastreável. \\
I5: fonte acadêmica, oficial ou comunitária tecnicamente verificável. & E5: é duplicata de outra fonte mais completa ou mais recente. \\
I6: disponível em inglês ou português. & E6: está abandonada ou inacessível para verificação. \\
I7: publicada ou atualizada entre 2018 e 2026. & E7: é apenas resumo, slide, anúncio breve ou opinião sem detalhe técnico. \\
I8: permite extrair pelo menos quatro campos do esquema de extração. & E8: menciona os frameworks apenas tangencialmente, sem análise substantiva. \\
\bottomrule
\end{tabularx}
\end{table}

Na literatura formal, a busca multibase retornou 100 registros brutos, reduzidos a 85 candidatos após a deduplicação por DOI e título, a aplicação do filtro temporal de 2018 a 2026 e a inclusão de uma semente dirigida pertencente ao escopo inicial e ausente da recuperação automática. Esses 85 candidatos passaram por uma primeira triagem de título e resumo, que incluiu 28 estudos para leitura de texto completo e excluiu 57. Na segunda passada, os 28 textos completos foram lidos integralmente sob os mesmos critérios, resultando em 25 estudos incluídos e 3 excluídos. As contagens e os motivos de exclusão de ambas as passadas foram registrados em log append-only para auditoria.

A partir dos 25 estudos incluídos, aplicou-se snowballing backward e forward segundo as diretrizes de Wohlin \cite{wohlin2014snowballing}, com um nível de profundidade e limite de três rodadas, cada uma semeada pelos incluídos da anterior. As três rodadas geraram, respectivamente, 69, 50 e 30 candidatos novos; destes, 9, 5 e 4 avançaram para texto completo e 6, 3 e 3 foram incluídos, e a busca forward não retornou citações elegíveis em nenhuma rodada. Ao todo, foram examinados 149 candidatos adicionais e incorporados 12 estudos novos, consolidando o corpus final em 37 estudos incluídos. Uma fonte aprovada na triagem da terceira rodada ficou sem texto completo na data de corte e foi mantida como pendência operacional, não como exclusão.

A revisão foi conduzida como multivocal por uma razão de cobertura: parte expressiva dos frameworks que organizam o desenvolvimento sobre um agente circula em fontes primárias não indexadas, como documentação oficial, repositórios e blogs de produto, de modo que restringir a revisão a artigos indexados deixaria essa camada de fora. A vertente de literatura cinzenta foi, portanto, levantada por busca dirigida a essas fontes primárias, sob o mesmo esquema de extração da literatura formal, mas com critérios próprios de escopo e tração. A caracterização aprofundada dos frameworks de apoio ao desenvolvimento por quem utiliza um agente, como GitHub Spec Kit, OpenSpec, BMAD Method, Get Shit Done e Spec Kitty, é reportada em um estudo companheiro dedicado a esse recorte; este artigo concentra-se no corpus formal de 37 estudos, que estabelece a estrutura e as tensões do campo. Reversa integra esse corpus formal por ter entrado como semente dirigida na busca.

De cada estudo incluído foram extraídos dez campos estruturados: nome, natureza, origem, artefatos, papéis, execução, validação, portabilidade, evidência e riscos. A consolidação da extração em matriz permitiu comparar os estudos de forma uniforme. A síntese seguiu abordagem temática em modo híbrido: temas indutivos foram identificados a partir das fichas de extração e, em seguida, confrontados com a taxonomia de seis dimensões proposta na revisão (especificação, contexto, papéis, execução, validação e portabilidade), reportando a sobreposição, as dimensões pouco cobertas e os temas que emergem fora da taxonomia.

Por fim, registre-se que a avaliação formal de qualidade por checklist não foi aplicada nesta rodada: a força da evidência é classificada por categoria de fonte, distinguindo estudo empírico, paper conceitual, documentação oficial e repositório público. Como repositórios e documentação de ferramentas mudam rapidamente, observações sobre versões, integrações, releases ou comandos devem ser lidas como um retrato do corpus consultado na data de corte, não como propriedades estáveis dos frameworks. As demais limitações que delimitam a leitura dos resultados são consolidadas na seção de Limitações da Revisão.
